
When a language model generates a TriviaQA answer, it leaves two independent traces of its confidence: how consistently its stochastic outputs cluster into the same semantic equivalence class (Semantic Entropy; SE), and how confidently it predicts each individual token on a deterministic greedy decode pass (minimum log-probability; min\_logprob). On TriviaQA dev with Llama-3.1-8B-Instruct across N=2500 prompts, the Pearson correlation between these two signals is |r| = 0.026 --- empirically near-zero, far below the collinearity threshold of 0.70 and the reparameterization boundary of 0.85. This near-orthogonality, while theoretically motivated by the algebraic structure of each signal, had not been directly measured before in the open-weight LLM hallucination detection literature.

Large language models are increasingly deployed in high-stakes domains --- healthcare, legal analysis, knowledge retrieval --- where hallucinated outputs carry real consequences. Uncertainty quantification (UQ) is the primary mechanism for flagging untrustworthy predictions, allowing downstream systems to abstain, escalate, or request human review. The need for accurate, calibrated UQ signals for open-weight LLMs has motivated a large and growing body of work (Guo et al., 2017; Xiong et al., 2023; Kang et al., 2025).

Two dominant signal families have emerged. First, token log-probability signals --- minimum log-probability, mean log-probability, sequence perplexity --- measure the model's token-level confidence from a single greedy decoding pass and require only standard generation API access \citep{Malinin2021Uncertainty}. On TriviaQA dev with Llama-3.1-8B, min\_logprob achieves AUROC approximately 0.825 in a prior internal pipeline experiment (h-m1, under identical experimental conditions --- not published externally). Second, consistency-based signals --- most notably Semantic Entropy \citep{Kuhn2023Semantic} --- measure the semantic diversity of N stochastic samples, capturing uncertainty at the semantic, sentence level rather than the token level. SE has shown strong standalone AUROC performance on factual QA benchmarks (Kuhn et al., 2023; Bouchard and Chauhan, 2025).

The natural next step --- combining SE and min\_logprob in an ensemble --- rests on a critical assumption: that these signals are statistically independent, meaning SE captures uncertainty dimensions not already captured by min\_logprob. Without measuring this independence, an ensemble might merely recombine redundant information. Remarkably, this assumption had not been directly tested. The closest related measurement --- first-token confidence versus SE correlation --- yields r = 0.54--0.76 on Llama-3.1-8B, TriviaQA \citep{Gabriel2026FirstToken}, suggesting moderate correlation for first-token signals. Whether min\_logprob (minimum over all tokens in greedy decode, not just the first) exhibits the same correlation with SE was an open question.

The key insight driving this work is that SE and min\_logprob are algebraically distinct operations by construction: SE marginalizes over semantic equivalence class probability masses computed from N=5 stochastic samples at temperature 0.7, while min\_logprob extracts the minimum token log-probability from a single greedy decoding path. These operations run on different inference calls, aggregate at different granularities (semantic cluster entropy vs. token-level minimum), and access different information about the model's output distribution. This algebraic separation manifests as statistical independence: Pearson |r| = 0.026 at N=2500.

This paper proceeds from an earlier proof-of-concept (PoC) run at N=300 (h-e1, results: Pearson |r|=0.049, partial $R^{2}=0.0101$, LRT p=0.1504) that validated the pipeline and confirmed the primary independence claim, but was underpowered for the partial $R^2$ threshold. The N=2500 evaluation reported throughout this paper supersedes the PoC and constitutes the primary experimental record.

\textbf{Contributions:}

\begin{enumerate}
\item \textbf{Empirical independence measurement:} Direct measurement of Pearson |r|(SE\_N5, min\_logprob) = 0.026 on TriviaQA dev with Llama-3.1-8B-Instruct (N=2500), confirming near-orthogonality and ruling out SE as a reparameterization of min\_logprob.
\end{enumerate}

\begin{enumerate}
\item \textbf{Circularity-controlled evaluation protocol:} Cross-model LM-judge design (Qwen-2.5-7B evaluating Llama-3.1-8B outputs) yields Spearman $\rho$(SE, judge) = $-0.026$, a replicable best practice for bias-robust UQ evaluation. This value shares the same rounded magnitude as the Pearson |r| above; both are computed on different variable pairs (SE vs. min\_logprob; SE vs. judge labels) via different correlation measures, and the coincident magnitude is a feature of this evaluation's near-zero correlation structure, not a transcription artifact.
\end{enumerate}

\begin{enumerate}
\item \textbf{Powered null result for linear conditional contribution:} At N=2500, SE adds no meaningful conditional predictive power beyond min\_logprob (partial $R^{2}=0.0005$, LRT p=0.442). The pre-registered gate evaluation classified this as \texttt{EXPLORE\textbackslash \_N10} (gate\_pass: false), indicating the pipeline recommends extending to N=10 stochastic samples. The partial $R^{2}=0.0005$ result is characterized as a powered null for the pre-registered effect size (estimated power >0.99 for detecting partial $R^{2}\geq 0.02$ at N=2500); the gate disclosure reflects the pipeline's conservative marginal-uncertainty criterion.
\end{enumerate}

\begin{enumerate}
\item \textbf{End-to-end reproducible pipeline:} A checkpoint-aware, GPU-efficient implementation completing the full N=2500 evaluation on TriviaQA dev, demonstrating production-quality execution for this hypothesis.
\end{enumerate}

